% GENERATED FILE. Edit canonical lessons, then run npm run generate:books.
% canonical-source-hash: fnv1a64:26444d8a9937a964
% canonical-lessons: GU-C340-uddes, GU-C340-hetu, GU-C340-irado, GU-C340-bahanu, GU-C340-sadhvu

\chapter{What It Is For}
\label{ch:gu-what-it-is-for}
\hlchaptermodality{340}

\begin{chapteropening}
\textbf{By the end of this chapter:} \emph{I can say what something is for and what someone means to do.}

Complete the last lesson of chapter 340: I can say what something is for and what someone means to do.
\end{chapteropening}

\section[\texorpdfstring{\gu{ઉદ્દેશ} (uddeś)}{ઉદ્દેશ (uddeś)}]{\gu{ઉદ્દેશ} (uddeś) — a purpose}

\label{lesson:GU-C340-uddes}



\begin{quote}
Before the new one: say the Gujarati for a situation, then the Gujarati for a circumstance.
\end{quote}

\subsection*{You'll want to know: \gu{ઉદ્દેશ}}
\textbf{\gu{ઉદ્દેશ}} — \emph{uddeś} — \textquotedblleft{}a purpose\textquotedblright{}.

\begin{quote}
\gu{આ} \gu{યોજનાનો} \gu{ઉદ્દેશ} \gu{ગામમાં} \gu{પાણી} \gu{લાવવાનો} \gu{છે}.
\end{quote}

\begin{quote}
\emph{ā yojnāno uddeś gāmmā̃ pāṇī lāvvāno che.}
\end{quote}

\begin{quote}
The aim of this plan is to bring water to the village.
\end{quote}

A masculine noun.

\begin{cousinweb}
\textbf{\gu{ઉદ્દેશ}} (\emph{uddeś}) is a tatsama, borrowed from Sanskrit \emph{uddeśa}, \textquotedblleft{}pointing at, an aim\textquotedblright{} (Monier-Williams).
\end{cousinweb}

\begin{grammarlens}[title={What you've built}]
One more word for giving reasons and opinions.
\end{grammarlens}

\subsection*{Your turn}
\begin{itemize}
\raggedright
  \item \emph{Say it:} \emph{uddeś}
  \item \emph{Say it:} \emph{uddeś}, once more
  \item \emph{Say it:} \emph{ā yojnāno uddeś gāmmā̃ pāṇī lāvvāno che.} — The aim of this plan is to bring water to the village.
  \item \emph{Recall:} say the Gujarati for a situation, then the Gujarati for a circumstance, then say \emph{uddeś} again
\end{itemize}

\begin{tcolorbox}[breakable,colback=teal!4,colframe=teal!35!black,title={Before you move on}]
What does \gu{ઉદ્દેશ} mean? (\textquotedblleft{}A purpose\textquotedblright{}.) Say it once more.
\end{tcolorbox}
\section[\texorpdfstring{\gu{હેતુ} (hetu)}{હેતુ (hetu)}]{\gu{હેતુ} (hetu) — a motive}

\label{lesson:GU-C340-hetu}



\begin{quote}
Before the new one: say the Gujarati for a circumstance, then the Gujarati for a purpose.
\end{quote}

\subsection*{You'll want to know: \gu{હેતુ}}
\textbf{\gu{હેતુ}} — \emph{hetu} — \textquotedblleft{}a motive\textquotedblright{}.

\begin{quote}
\gu{તેના} \gu{આવવાનો} \gu{હેતુ} \gu{શું} \gu{હતો}?
\end{quote}

\begin{quote}
\emph{tenā āvvāno hetu śũ hato?}
\end{quote}

\begin{quote}
What was his motive in coming?
\end{quote}

A masculine noun.

\begin{cousinweb}
\textbf{\gu{હેતુ}} (\emph{hetu}) is a tatsama, borrowed from Sanskrit \emph{hetu}, \textquotedblleft{}a cause, a motive\textquotedblright{} (Monier-Williams).
\end{cousinweb}

\begin{grammarlens}[title={What you've built}]
One more word for giving reasons and opinions.
\end{grammarlens}

\subsection*{Your turn}
\begin{itemize}
\raggedright
  \item \emph{Say it:} \emph{hetu}
  \item \emph{Say it:} \emph{hetu}, once more
  \item \emph{Say it:} \emph{tenā āvvāno hetu śũ hato?} — What was his motive in coming?
  \item \emph{Recall:} say the Gujarati for a circumstance, then the Gujarati for a purpose, then say \emph{hetu} again
\end{itemize}

\begin{tcolorbox}[breakable,colback=teal!4,colframe=teal!35!black,title={Before you move on}]
What does \gu{હેતુ} mean? (\textquotedblleft{}A motive\textquotedblright{}.) Say it once more.
\end{tcolorbox}
\section[\texorpdfstring{\gu{ઇરાદો} (irādo)}{ઇરાદો (irādo)}]{\gu{ઇરાદો} (irādo) — an intention}

\label{lesson:GU-C340-irado}



\begin{quote}
Before the new one: say the Gujarati for a purpose, then the Gujarati for a motive.
\end{quote}

\subsection*{You'll want to know: \gu{ઇરાદો}}
\textbf{\gu{ઇરાદો}} — \emph{irādo} — \textquotedblleft{}an intention\textquotedblright{}.

\begin{quote}
\gu{અરુણનો} \gu{ઇરાદો} \gu{સારો} \gu{હતો}, \gu{પણ} \gu{પરિણામ} \gu{ખરાબ} \gu{આવ્યું}.
\end{quote}

\begin{quote}
\emph{Aruṇno irādo sāro hato, paṇ pariṇām kharāb āvyũ.}
\end{quote}

\begin{quote}
Arun meant well, but the result was bad.
\end{quote}

A masculine noun.

\begin{cousinweb}
\textbf{\gu{ઇરાદો}} (\emph{irādo}) is usually traced to Arabic \emph{irāda}, \textquotedblleft{}will, intention\textquotedblright{}, which came to Gujarati through Persian.
\end{cousinweb}

\begin{grammarlens}[title={What you've built}]
One more word for giving reasons and opinions.
\end{grammarlens}

\subsection*{Your turn}
\begin{itemize}
\raggedright
  \item \emph{Say it:} \emph{irādo}
  \item \emph{Say it:} \emph{irādo}, once more
  \item \emph{Say it:} \emph{Aruṇno irādo sāro hato, paṇ pariṇām kharāb āvyũ.} — Arun meant well, but the result was bad.
  \item \emph{Recall:} say the Gujarati for a purpose, then the Gujarati for a motive, then say \emph{irādo} again
\end{itemize}

\begin{tcolorbox}[breakable,colback=teal!4,colframe=teal!35!black,title={Before you move on}]
What does \gu{ઇરાદો} mean? (\textquotedblleft{}An intention\textquotedblright{}.) Say it once more.
\end{tcolorbox}
\section[\texorpdfstring{\gu{બહાનું} (bahānũ)}{બહાનું (bahānũ)}]{\gu{બહાનું} (bahānũ) — an excuse}

\label{lesson:GU-C340-bahanu}



\begin{quote}
Before the new one: say the Gujarati for a motive, then the Gujarati for an intention.
\end{quote}

\subsection*{You'll want to know: \gu{બહાનું}}
\textbf{\gu{બહાનું}} — \emph{bahānũ} — \textquotedblleft{}an excuse\textquotedblright{}.

\begin{quote}
\gu{તેણે} \gu{મોડા} \gu{આવવાનું} \gu{એક} \gu{નવું} \gu{બહાનું} \gu{કાઢ્યું}.
\end{quote}

\begin{quote}
\emph{teṇe moḍā āvvānũ ek navũ bahānũ kāḍhyũ.}
\end{quote}

\begin{quote}
He came up with a new excuse for being late.
\end{quote}

A neuter noun.

\begin{cousinweb}
\textbf{\gu{બહાનું}} (\emph{bahānũ}) is usually traced to Persian \emph{bahāna}, \textquotedblleft{}a pretext\textquotedblright{}.
\end{cousinweb}

\begin{grammarlens}[title={What you've built}]
One more word for giving reasons and opinions.
\end{grammarlens}

\subsection*{Your turn}
\begin{itemize}
\raggedright
  \item \emph{Say it:} \emph{bahānũ}
  \item \emph{Say it:} \emph{bahānũ}, once more
  \item \emph{Say it:} \emph{teṇe moḍā āvvānũ ek navũ bahānũ kāḍhyũ.} — He came up with a new excuse for being late.
  \item \emph{Recall:} say the Gujarati for a motive, then the Gujarati for an intention, then say \emph{bahānũ} again
\end{itemize}

\begin{tcolorbox}[breakable,colback=teal!4,colframe=teal!35!black,title={Before you move on}]
What does \gu{બહાનું} mean? (\textquotedblleft{}An excuse\textquotedblright{}.) Say it once more.
\end{tcolorbox}
\section[\texorpdfstring{\gu{સાધવું} (sādhvũ)}{સાધવું (sādhvũ)}]{\gu{સાધવું} (sādhvũ) — to accomplish}

\label{lesson:GU-C340-sadhvu}



\begin{quote}
Before the new one: say the Gujarati for an intention, then the Gujarati for an excuse.
\end{quote}

\subsection*{You'll want to know: \gu{સાધવું}}
\textbf{\gu{સાધવું}} — \emph{sādhvũ} — \textquotedblleft{}to accomplish\textquotedblright{}.

\begin{quote}
\gu{મહેનતથી} \gu{તે} \gu{પોતાનો} \gu{હેતુ} \gu{સાધશે}.
\end{quote}

\begin{quote}
\emph{mahenatthī te potāno hetu sādhśe.}
\end{quote}

\begin{quote}
With hard work he will achieve his aim.
\end{quote}

Past, with \gu{તેણે} (he or she …): \gu{તેણે} \gu{પોતાનો} \gu{હેતુ} \gu{સાધ્યો}, \emph{teṇe potāno hetu sādhyo}.

\begin{cousinweb}
\textbf{\gu{સાધવું}} (\emph{sādhvũ}) is inherited from Sanskrit \emph{sādhnōti}, \textquotedblleft{}accomplishes\textquotedblright{}; CDIAL lists the Gujarati verb (CDIAL 13339).
\end{cousinweb}

\begin{grammarlens}[title={What you've built}]
One more word for giving reasons and opinions.
\end{grammarlens}

\subsection*{Your turn}
\begin{itemize}
\raggedright
  \item \emph{Say it:} \emph{sādhvũ}
  \item \emph{Say it:} \emph{sādhvũ}, once more
  \item \emph{Say it:} \emph{mahenatthī te potāno hetu sādhśe.} — With hard work he will achieve his aim.
  \item \emph{Recall:} say the Gujarati for an intention, then the Gujarati for an excuse, then say \emph{sādhvũ} again
\end{itemize}

\begin{tcolorbox}[breakable,colback=teal!4,colframe=teal!35!black,title={Before you move on}]
What does \gu{સાધવું} mean? (\textquotedblleft{}To accomplish\textquotedblright{}.) Say it once more.
\end{tcolorbox}
